% ch03_u01 —— 原书 p.61–p.70（3.1 全部、3.2 全部、3.3.1）
\SourcePage{61}

\bdsec{3.1. 定义与基本性质}

\bdp{3.1.1} 设 $G$ 是 $\mathbb{C}$ 上的连通约化群，其上固定 Borel 子群 $B = B_G \subset G$。令 $N = [B,B]$，于是 $H = B/N$ 是嘉当群。记相应的李代数为 $\mathfrak{n} \subset \mathfrak{b} \subset \mathfrak{g}$，$\mathfrak{h} = \mathfrak{b}/\mathfrak{n}$。$\mathfrak{g}$ 带一个典范的降李代数滤链 $\mathfrak{g}^k$，使得 $\mathfrak{g}^0 = \mathfrak{b}$，$\mathfrak{g}^1 = \mathfrak{n}$；且对任意 $k > 0$，$\mathfrak{h} = \operatorname{gr}^0\mathfrak{g}$ 在 $\operatorname{gr}^k\mathfrak{g}$（相应地 $\operatorname{gr}^{-k}\mathfrak{g}$）上作用的权，都是 $k$ 个单正根（相应地单负根）之和。特别地 $\operatorname{gr}^{-1}\mathfrak{g} = \bigoplus \mathfrak{g}^{\alpha}$，其中 $\alpha$ 是单负根。令 $Z = Z_G = \operatorname{Center} G$。

\bdp{3.1.2} 设 $X$ 是任意光滑曲线（不必完备），$\mathfrak{F}_B$ 是其上的 $B$-丛。记 $\mathfrak{F}_G$ 为诱导出的 $G$-torsor，于是 $\mathfrak{F}_B \subset \mathfrak{F}_G$。有相应的扭李代数 $\mathfrak{b}_{\mathfrak{F}} := \mathfrak{b}_{\mathfrak{F}_B}$ 与 $\mathfrak{g}_{\mathfrak{F}} := \mathfrak{g}_{\mathfrak{F}_B} = \mathfrak{g}_{\mathfrak{F}_G}$，它们带李代数滤链 $\mathfrak{g}^k_{\mathfrak{F}}$。考虑联络层 $\operatorname{Conn}(\mathfrak{F}_B)$、$\operatorname{Conn}(\mathfrak{F}_G)$：它们分别是 $\mathfrak{b}_{\mathfrak{F}}\otimes\omega_X$- 与 $\mathfrak{g}_{\mathfrak{F}}\otimes\omega_X$-torsor（主齐性空间）。有显然的嵌入 $\operatorname{Conn}(\mathfrak{F}_B) \subset \operatorname{Conn}(\mathfrak{F}_G)$。它定义投影 $c: \operatorname{Conn}(\mathfrak{F}_G) \to (\mathfrak{g}/\mathfrak{b})_{\mathfrak{F}}\otimes\omega_X$，使得 $c^{-1}(0) = \operatorname{Conn}(\mathfrak{F}_B)$，且对任意 $\nabla \in \operatorname{Conn}(\mathfrak{F}_G)$ 与 $\nu \in \mathfrak{g}_{\mathfrak{F}}\otimes\omega_X$ 有 $c(\nabla+\nu) = c(\nabla)+\nu \bmod \mathfrak{b}_{\mathfrak{F}}\otimes\omega_X$。

\bdp{3.1.3} \emph{定义.} $X$ 上的 $G$-oper 是一对 $(\mathfrak{F}_B,\nabla)$，其中 $\nabla \in \Gamma(X,\operatorname{Conn}(\mathfrak{F}_G))$，且满足

\begin{enumerate}[leftmargin=2.2em]
\item[(1)] $c(\nabla) \in \operatorname{gr}^{-1}\mathfrak{g}_{\mathfrak{F}}\otimes\omega_X \subset (\mathfrak{g}/\mathfrak{b})_{\mathfrak{F}}\otimes\omega_X$；
\item[(2)] 对任意单负根 $\alpha$，$\alpha$-分量 $c(\nabla)^{\alpha} \in \Gamma(X,\mathfrak{g}^{\alpha}_{\mathfrak{F}}\otimes\omega_X)$ 在 $X$ 的任一点都不为零。
\end{enumerate}

若 $\mathfrak{g}$ 是半单李代数，则 $\mathfrak{g}$-oper 就是 $G_{ad}$-oper，其中 $G_{ad}$ 是对应于 $\mathfrak{g}$ 的伴随群。

我们通常把 $G$-oper 看作 $G$-局部系统 $(\mathfrak{F}_G,\nabla)$，再附加 oper 结构，即满足上述条件 (1)、(2) 的 $B$-旗 $\mathfrak{F}_B \subset \mathfrak{F}_G$。

\SourcePage{62}

$X$ 上的 $G$-oper 构成群胚 $\operatorname{Op}_G(X)$。当 $X'$ 在 $X$ 上平展时，群胚 $\operatorname{Op}_G(X')$ 构成 $X_{\text{\'et}}$ 上的群胚层 $\operatorname{Op}_G$。

\bdp{3.1.4} \emph{命题.} 设 $(\mathfrak{F}_B,\nabla)$ 是 $G$-oper。若 $X$ 连通，则 $\operatorname{Aut}(\mathfrak{F}_B,\nabla) = Z$。\hfill$\Box$

特别地，$\mathfrak{g}$-oper 没有对称性，即 $\operatorname{Op}_{\mathfrak{g}}(X)$ 是集合，$\operatorname{Op}_{\mathfrak{g}}$ 是集合层。

\bdp{3.1.5} \emph{命题.} 设 $X$ 完备、连通，亏格 $g > 1$。设 $(\mathfrak{F}_G,\nabla)$ 是 $X$ 上的 $G$-局部系统，且具有 oper 结构。则

\begin{enumerate}[label=(\roman*),leftmargin=2.2em]
\item $(\mathfrak{F}_G,\nabla)$ 上的 oper 结构唯一：相应的旗 $\mathfrak{F}_B \subset \mathfrak{F}_G$ 是 Harder--Narasimhan 旗；
\item $\operatorname{Aut}(\mathfrak{F}_G,\nabla) = Z$；
\item $(\mathfrak{F}_G,\nabla)$ 不能约化到非平凡抛物子群 $P \subset G$。\hfill$\Box$
\end{enumerate}

当然，ii) 由 i) 与 3.1.4 推出。

\bdp{3.1.6} \emph{例.} $GL_n$-oper 可以看作 $\mathcal{O}_X$-模 $\mathcal{E}$，带联络 $\nabla: \mathcal{E} \to \mathcal{E}\otimes\omega_X$ 与滤链 $\mathcal{E} = \mathcal{E}_n \supset \mathcal{E}_{n-1} \supset \cdots \supset \mathcal{E}_0 = 0$，使得

\begin{enumerate}[label=(\roman*),leftmargin=2.2em]
\item 层 $\operatorname{gr}_i\mathcal{E}$（$n \ge i \ge 1$）可逆；
\item $\nabla(\mathcal{E}_i) \subset \mathcal{E}_{i+1}\otimes\omega_X$，且对 $n-1 \ge i \ge 1$，由 $\nabla$ 诱导的态射 $\operatorname{gr}_i\mathcal{E} \to \operatorname{gr}_{i+1}\mathcal{E}\otimes\omega_X$ 是同构。
\end{enumerate}

$GL_n$-oper 可以如下构造。设 $\mathcal{A}$、$\mathcal{B}$ 是可逆 $\mathcal{O}_X$-模，$\partial: \mathcal{A} \to \mathcal{B}$ 是 $n$ 阶微分算子，其象征 $\sigma(\partial) \in \Gamma(X,\mathcal{B}\otimes\mathcal{A}^{\otimes(-1)}\otimes\Theta_X^{\otimes n})$ 无处为零。这里的 $\partial$ 是 $\mathcal{B}\otimes_{\mathcal{O}_X}\mathcal{D}_X\otimes_{\mathcal{O}_X}\mathcal{A}^{\otimes(-1)}$ 的截面；等价地，是 $\mathcal{O}$-线性映射 $\mathcal{B}^{\otimes(-1)} \to \mathcal{D}_X\otimes\mathcal{A}^{\otimes(-1)}$。设 $I \subset \mathcal{D}_X\otimes\mathcal{A}^{\otimes(-1)}$ 是该映射的像生成的 $\mathcal{D}_X$-子模。令 $\mathcal{E} := \mathcal{D}_X\otimes\mathcal{A}^{\otimes(-1)}/I$；记 $\mathcal{E}_i$ 为 $\mathcal{D}_X$ 上按算子阶数给出的通常滤链在 $\mathcal{E}$ 上诱导的滤链。\SourcePage{63}

于是 $\mathcal{E}$ 是 $\mathcal{D}_X$-模，即带联络的 $\mathcal{O}_X$-模；滤链 $\mathcal{E}_i$ 满足条件 (i)、(ii)。因此 $(\mathcal{E},\{\mathcal{E}_i\},\nabla)$ 是 $GL_n$-oper。这一构造给出 $GL_n$-oper 的群胚与上述数据 $\partial: \mathcal{A} \to \mathcal{B}$ 的群胚之间的等价。

逆函子 $\Phi$ 把 $(\mathcal{E},\{\mathcal{E}_i\},\nabla)$ 对应到如下微分算子：$\mathcal{A} \to \mathcal{B}$，其中 $\mathcal{A} := \mathcal{E}_1^{\otimes(-1)}$，$\mathcal{B} := \omega_X\otimes(\mathcal{E}/\mathcal{E}_{n-1})^{\otimes(-1)}$。把 $\mathcal{E}$ 看作 $\mathcal{D}_X$-模。设 $\mathcal{D}_X^{(k)} \subset \mathcal{D}_X$ 是阶数 $\le k$ 的算子构成的子层。则态射 $\mathcal{D}_X^{(n-1)}\otimes_{\mathcal{O}_X}\mathcal{E}_1 \to \mathcal{E}$ 是同构；因此复合 $\mathcal{D}_X^{(n)}\otimes_{\mathcal{O}_X}\mathcal{E}_1 \to \mathcal{E} \xrightarrow{\sim} \mathcal{D}_X^{(n-1)}\otimes_{\mathcal{O}_X}\mathcal{E}_1$ 给出正合列 $0 \to \mathcal{D}_X^{(n-1)}\otimes_{\mathcal{O}_X}\mathcal{E}_1 \to \mathcal{D}_X^{(n)}\otimes_{\mathcal{O}_X}\mathcal{E}_1 \to \omega_X^{\otimes(-n)}\otimes\mathcal{E}_1 \to 0$ 的分裂，即一个态射 $\omega_X^{\otimes(-n)}\otimes\mathcal{E}_1 \to \mathcal{D}_X^{(n)}\otimes_{\mathcal{O}_X}\mathcal{E}_1$。这等同于一个微分算子 $\partial: \mathcal{A} \to \mathcal{B}$（注意：同构 $\operatorname{gr}_i\mathcal{E} \xrightarrow{\sim} \operatorname{gr}_{i+1}\mathcal{E}\otimes\omega_X$ 诱导同构 $\mathcal{E}_1 \xrightarrow{\sim} (\mathcal{E}/\mathcal{E}_{n-1})\otimes\omega_X^{\otimes(n-1)}$，故 $\omega_X^{\otimes(-n)}\otimes\mathcal{E}_1 = \omega_X^{\otimes(-1)}\otimes(\mathcal{E}/\mathcal{E}_{n-1}) = \mathcal{B}^{\otimes(-1)}$）。

把上述函子 $\Phi$ 施于 $\mathrm{SL}_2$-oper，得到微分算子 $\partial: \mathcal{A} \to \omega_X\otimes\mathcal{A}^{\otimes(-1)}$。容易证明，这样得到 $\mathrm{SL}_2$-oper 的群胚与下述对象的群胚之间的等价：由可逆层 $\mathcal{A}$ 与 Sturm--Liouville 算子 $\partial: \mathcal{A} \to \omega_X\otimes\mathcal{A}^{\otimes(-1)}$ 构成的对 $(\mathcal{A},\partial)$，即符号 $\sigma(\partial)$ 无处为零的自伴二阶微分算子 $\partial$。注意 $\sigma(\partial)$ 诱导同构 $\omega_X^{\otimes 2}\otimes\mathcal{A} \xrightarrow{\sim} \omega_X\otimes\mathcal{A}^{\otimes(-1)}$，所以 $\mathcal{A}$ 自动是 $\omega_X^{\otimes(-1)}$ 的平方根。

若 $(\mathcal{A},\partial)$ 是 Sturm--Liouville 算子，$\mathcal{M}$ 是带同构 $\mathcal{M}^{\otimes 2} \xrightarrow{\sim} \mathcal{O}_X$ 的线丛，则 $\mathcal{M}$ 有典范联络；于是用 $\mathcal{M}$ 扭 $(\mathcal{A},\partial)$，得到 Sturm--Liouville 算子 $(\tilde{\mathcal{A}},\tilde{\partial})$，$\tilde{\mathcal{A}} = \mathcal{A}\otimes\mathcal{M}$。称 $(\mathcal{A},\partial)$ 与 $(\tilde{\mathcal{A}},\tilde{\partial})$ 等价。容易看出，自然映射 $\operatorname{Op}_{\mathrm{SL}_2}(X) \to \operatorname{Op}_{\mathfrak{sl}_2}(X)$ 把 $\operatorname{Op}_{\mathfrak{sl}_2}(X)$ 与 Sturm--Liouville 算子的等价类集合等同。

其它典型群的 oper 可用类似方式描述（局部情形已在 [DS85, 第 8 节] 中完成）。

\SourcePage{64}

\bdp{3.1.7} 把 $\mathfrak{sl}_2$-oper 与 Sturm--Liouville 算子的等价类等同（见 3.1.6），可见 $\operatorname{Op}_{\mathfrak{sl}_2}$ 是 $\omega_X^{\otimes 2}$-torsor：$\omega_X^{\otimes 2}$ 的截面 $\eta$ 把 Sturm--Liouville 算子 $\partial: \mathcal{A} \to \mathcal{A}\otimes\omega_X^{\otimes 2}$（$\mathcal{A}^{\otimes(-2)} = \omega_X$）映到 $\partial - \eta$。下面不用 Sturm--Liouville 算子，来描述 $\omega_X^{\otimes 2}$ 在 $\operatorname{Op}_{\mathfrak{sl}_2}$ 上的这个作用。

利用 $\mathfrak{sl}_2$ 上的双线性型 $\operatorname{Tr}(AB)$，把 $\mathfrak{n} \subset \mathfrak{sl}_2$ 与 $(\mathfrak{sl}_2/\mathfrak{b})^*$ 等同。若 $\mathfrak{F} = (\mathfrak{F}_B,\nabla)$ 是 $\mathfrak{sl}_2$-oper，则由 3.1.3，截面 $c(\nabla)$ 平凡化了层 $(\mathfrak{sl}_2/\mathfrak{b})_{\mathfrak{F}_B}\otimes\omega_X$。于是 $(\mathfrak{sl}_2/\mathfrak{b})_{\mathfrak{F}_B} = \omega_X^{\otimes(-1)}$，$\mathfrak{n}_{\mathfrak{F}} = \omega_X$，并有嵌入 $\omega_X^{\otimes 2} = \mathfrak{n}_{\mathfrak{F}_B}\otimes\omega_X \hookrightarrow (\mathfrak{sl}_2)_{\mathfrak{F}_B}\otimes\omega_X$。用 $\omega_X^{\otimes 2} \subset (\mathfrak{sl}_2)_{\mathfrak{F}_B}\otimes\omega_X$ 的截面 $\mu$ 平移 $\nabla$，得到新的 oper，记作 $\mathfrak{F}+\mu$。$\operatorname{Op}_{\mathfrak{sl}_2}$ 上的这个 $\omega_X^{\otimes 2}$-作用与上面引入的作用一致，因此它使 $\operatorname{Op}_{\mathfrak{sl}_2}$ 成为 $\omega_X^{\otimes 2}$-torsor。

\emph{注.} 众所周知，这个 torsor 平凡（即使 $H^1(X,\omega_X^{\otimes 2}) \ne 0$，即 $g \le 1$；$\mathbb{P}^1$ 上或椭圆曲线上确实存在 Sturm--Liouville 算子）。然而对曲线的族 $X$，这个 torsor 可能不平凡。

\bdp{3.1.8} 3.1.9 中要用如下记号。设 $B_0 \subset PSL_2$ 是上三角矩阵构成的群。令 $N_0 := [B_0,B_0]$，$\mathfrak{b}_0 := \operatorname{Lie} B_0$，$\mathfrak{n}_0 := \operatorname{Lie} N_0$。借助伴随作用 $B_0/N_0 \to \operatorname{Aut}\mathfrak{n}_0 = \mathbb{G}_m$，把 $B_0/N_0$ 与 $\mathbb{G}_m$ 等同。用矩阵 $e := \begin{pmatrix}0&1\\0&0\end{pmatrix}$ 与 $f := \begin{pmatrix}0&0\\1&0\end{pmatrix}$，把 $\mathfrak{n}_0$ 与 $\mathfrak{sl}_2/\mathfrak{b}_0$ 都等同于 $\mathbb{C}$。于是对 $\mathfrak{sl}_2$-oper $\mathfrak{F} = (\mathfrak{F}_{B_0},\nabla)$，3.1.7 中的同构 $(\mathfrak{sl}_2/\mathfrak{b}_0)_{\mathfrak{F}_{B_0}} \xrightarrow{\sim} \omega_X^{\otimes(-1)}$（或同构 $\mathfrak{n}_{\mathfrak{F}_{B_0}} \xrightarrow{\sim} \omega_X$）诱导同构：$\mathfrak{F}_{B_0}$ 沿 $B_0 \to B_0/N_0 = \mathbb{G}_m$ 的前推同构于 $\mathbb{G}_m$-torsor $\omega_X$。

\bdp{3.1.9} 对任意半单李代数 $\mathfrak{g}$，我们将给出 $\operatorname{Op}_{\mathfrak{g}}(X)$ 相当显式的描述。特别地，我们将在 $\operatorname{Op}_{\mathfrak{g}}(X)$ 上引入一个「典范的」仿射空间结构（$\mathfrak{g} = \mathfrak{sl}_2$ 时的情形已在 3.1.7 中引入）。设 $G$ 是对应于 $\mathfrak{g}$ 的伴随群，$B$ 是其 Borel 子群。沿用 3.1.8 的记号。固定主嵌入 $i: \mathfrak{sl}_2 \hookrightarrow \mathfrak{g}$，使 $i(\mathfrak{b}_0) \subset \mathfrak{b}$；有相应的嵌入 $i_G: PSL_2 \hookrightarrow G$，\SourcePage{65}$i_B: B_0 \hookrightarrow B$。

令 $V = V_{\mathfrak{g}} := \mathfrak{g}^{N_0}$。于是 $\mathfrak{n}_0 \subset V \subset \mathfrak{n}$。有 $\mathbb{G}_m = B_0/N_0$ 在 $V$ 上的伴随作用 $\operatorname{Ad}$。定义 $V$ 上新的 $\mathbb{G}_m$-作用 $a$：$a(t)v := t\operatorname{Ad}(t)v$，$v \in V$，$t \in \mathbb{G}_m$。

考虑向量丛 $V_{\omega_X}$，即 $V$ 关于 $\mathbb{G}_m$-作用 $a$ 的 $\omega_X$-扭（这里把 $\omega_X$ 看作 $X$ 上的 $\mathbb{G}_m$-torsor）。用 $\omega_X$ 扭嵌入 $\mathbb{C} \xrightarrow{\sim} \mathbb{C}e = \mathfrak{n}_0 \hookrightarrow V$，得到嵌入 $\omega_X^{\otimes 2} \hookrightarrow V_{\omega_X}$。

对任意 $\mathfrak{sl}_2$-oper $\mathfrak{F}_0 = (\mathfrak{F}_{B_0},\nabla_0)$，其 $i$-前推 $i\mathfrak{F}_0 = (\mathfrak{F}_B,\nabla)$ 是 $\mathfrak{g}$-oper。由 3.1.8 有典范同构 $V_{\omega_X} = V_{\mathfrak{F}_{B_0}}\otimes\omega_X$，从而有典范嵌入 $V_{\omega_X} \subset \mathfrak{b}_{\mathfrak{F}_0}\otimes\omega_X = \mathfrak{b}_{\mathfrak{F}_B}\otimes\omega_X$。用 $V_{\omega_X}$ 的截面 $\nu$ 平移 $\nabla$，得到新的 $\mathfrak{g}$-oper，记作 $i\mathfrak{F}_0+\nu$。

令 $\underline{\operatorname{Op}}_{\mathfrak{g}}$ 为 $\omega_X^{\otimes 2}$-torsor $\operatorname{Op}_{\mathfrak{sl}_2}$ 经嵌入 $\omega_X^{\otimes 2} \subset V_{\omega_X}$ 诱导出的 $V_{\omega_X}$-torsor。$\underline{\operatorname{Op}}_{\mathfrak{g}}$ 的截面是一对如上所述的 $(\mathfrak{F}_0,\nu)$；并约定：对 $\omega_X^{\otimes 2}$ 的截面 $\mu$，有 $(\mathfrak{F}_0+\mu,\nu) = (\mathfrak{F}_0,\mu+\nu)$。有典范映射
\begin{equation}\tag{43}
\underline{\operatorname{Op}}_{\mathfrak{g}} \longrightarrow \operatorname{Op}_{\mathfrak{g}},
\end{equation}
它把 $(\mathfrak{F}_0,\nu)$ 映到 $i\mathfrak{F}_0+\nu$。

\bdp{3.1.10} \emph{命题.} 映射（43）是双射。\hfill$\Box$

\emph{注.}
\begin{enumerate}[label=(\roman*),leftmargin=2.2em]
\item 尽管双射（43）是典范的，我们不确定它给出 $\operatorname{Op}_{\mathfrak{g}}$ 的合理描述。
\item 3.1.9 中的空间 $V = V_{\mathfrak{g}}$ 依赖于主嵌入 $i: \mathfrak{sl}_2 \hookrightarrow \mathfrak{g}$ 的选取（对这种 $i$，存在唯一的 Borel 子代数 $\mathfrak{b} \subset \mathfrak{g}$ 包含 $i(\mathfrak{b}_0)$）。但任意两个主嵌入 $\mathfrak{sl}_2 \hookrightarrow \mathfrak{g}$ 都由 $G = G_{ad}$ 的唯一元素共轭。于是可以等同对应于不同 $i$ 的各个 $V$，得到与 $\mathfrak{g}$ 典范相伴的向量空间（不是 $\mathfrak{g}$ 的子空间！）。
\item 设 $G$ 是对应于 $\mathfrak{g}$ 的伴随群，$B$ 是 $G$ 的 Borel 子群。命题 3.1.10 表明，对任意 $\mathfrak{g}$-oper $\mathfrak{F} = (\mathfrak{F}_B,\nabla)$，$\mathfrak{F}_B$ 同构于某个典范 $B$-丛 $\mathfrak{F}_B^0$，它不依赖于 $\mathfrak{F}$。\SourcePage{66}事实上，$\mathfrak{F}_B^0$ 是 2.6.5 中典范 $(\operatorname{Aut}^0 O)$-丛沿某个同态 $i_B\sigma\pi: \operatorname{Aut}^0 O \to B$ 的前推。这里 $\pi$ 是投影 $\operatorname{Aut}^0 O \to \operatorname{Aut}(O/\mathfrak{m}^3)$，$\mathfrak{m}$ 是 $O$ 的极大理想；$\sigma$ 是同构 $\operatorname{Aut}(O/\mathfrak{m}^3) \xrightarrow{\sim} B_0$，$B_0$ 是 $PSL_2$ 的 Borel 子群；$i_B: B_0 \to B$ 由主嵌入 $PSL_2 \to G$ 诱导（$\sigma$ 与 $i_B$ 在唯一共轭的意义下唯一）。
\end{enumerate}

\bdp{3.1.11} 设 $X$ 完备。则 $G$-oper 构成光滑代数栈，仍记作 $\operatorname{Op}_G(X)$（记号滥用）。若 $G$ 半单，这是 Deligne--Mumford 栈（见 3.1.4）；若 $G$ 伴随，则 $\operatorname{Op}_G(X) = \operatorname{Op}_{\mathfrak{g}}(X)$ 是概形，经（43）同构于仿射空间 $\underline{\operatorname{Op}}_{\mathfrak{g}}(X)$。

\emph{注.}
\begin{enumerate}[label=(\roman*),leftmargin=2.2em]
\item 若 $X$ 非完备，则 $\operatorname{Op}_{\mathfrak{g}}(X)$ 是归纳概形。
\item 若 $X$ 完备、连通，亏格 $g > 1$，则 $\dim\operatorname{Op}_{\mathfrak{g}}(X) = (g-1)\cdot\dim\mathfrak{g}$。事实上，由命题 3.1.10，$\dim\operatorname{Op}_{\mathfrak{g}}(X) = \dim\underline{\operatorname{Op}}_{\mathfrak{g}}(X) = \dim\Gamma(X,V_{\omega_X})$；Hitchin 的一个简单计算（见 2.2.4 的注 4）表明，当 $g > 1$ 时 $\dim\Gamma(X,V_{\omega_X}) = (g-1)\cdot\dim\mathfrak{g}$。实际上，3.1.13 将表明 $\Gamma(X,V_{\omega_X}) = \operatorname{Hitch}_{{}^{L}\mathfrak{g}}(X)$，所以直接使用 2.2.4(ii) 中提到的 Hitchin 公式即可：
\[
\dim\operatorname{Hitch}_{{}^{L}\mathfrak{g}}(X) = (g-1)\cdot\dim{}^{L}\mathfrak{g} = (g-1)\cdot\dim\mathfrak{g}.
\]
\item 设 $X$ 如注 ii，$G$ 是对应于 $\mathfrak{g}$ 的伴随群。有显然的态射 $i: \operatorname{Op}_{\mathfrak{g}}(X) \to \operatorname{LocSys}_G$，其中 $\operatorname{LocSys}_G$ 是 $X$ 上 $G$-局部系统的栈。可以证明：不能约化到非平凡抛物子群 $P \subset G$ 且没有非平凡自同构的 $G$-局部系统，构成开子栈 $U \subset \operatorname{LocSys}_G$。它实际上是光滑簇；$U$ 带典范辛结构。由 3.1.5，\SourcePage{67}$i(\operatorname{Op}_{\mathfrak{g}}(X)) \subset U$，且 $i$ 是集合论意义下的嵌入。事实上 $i$ 是闭嵌入，$i(\operatorname{Op}_{\mathfrak{g}}(X))$ 是 $U$ 的 Lagrange 子簇。此外，$i(\operatorname{Op}_{\mathfrak{g}}(X)) = \pi^{-1}(S)$，其中 $\pi: \operatorname{LocSys}_G \to \operatorname{Bun}_G$ 对应于忘掉联络，$S \subset \operatorname{Bun}_G$ 是同构于 $\mathfrak{F}_G^0$ 的 $G$-丛构成的局部闭子栈；$\mathfrak{F}_G^0$ 是对应于 3.1.10 注 iii 中引入的 $B$-丛 $\mathfrak{F}_B^0$ 的 $G$-丛（于是 $S$ 是幂幺群 $\operatorname{Aut}\mathfrak{F}_G^0$ 的分类栈）。
\end{enumerate}

\bdp{3.1.12} 记 $A_{\mathfrak{g}}(X)$ 为 $\operatorname{Op}_{\mathfrak{g}}(X)$ 的坐标环。我们将构造 $A_{\mathfrak{g}}(X)$ 上的典范滤链，以及分次代数的典范同构
\begin{equation}\tag{44}
\sigma_{A(X)}: \operatorname{gr} A_{\mathfrak{g}}(X) \xrightarrow{\sim} \mathfrak{z}^{cl}_{{}^{L}\mathfrak{g}}(X),
\end{equation}
其中 ${}^{L}\mathfrak{g}$ 表示 $\mathfrak{g}$ 的朗兰兹对偶，（44）的右端已在 2.2.2 中定义。我们给出两种等价的构造：3.1.13 中的构造直截了当，用到同构（43）；3.1.14 中的构造更自然。

\bdp{3.1.13} 利用 3.1.8，把 $A_{\mathfrak{g}}(X)$ 与 $\underline{\operatorname{Op}}_{\mathfrak{g}}(X)$ 的坐标环等同。记 $A^{cl}_{\mathfrak{g}}(X)$ 为向量空间 $\Gamma(X,V_{\omega_X})$ 的坐标环，它对应于仿射空间 $\underline{\operatorname{Op}}_{\mathfrak{g}}(X)$。考虑 $A^{cl}_{\mathfrak{g}}(X)$ 上与 $V$ 上 $\mathbb{G}_m$-作用 $a$ 所诱导者相反的作用（见 3.1.7）；$A^{cl}_{\mathfrak{g}}(X)$ 上相应的分次为正。它诱导 $A_{\mathfrak{g}}(X)$ 上的典范环滤链，以及典范同构 $\operatorname{gr} A_{\mathfrak{g}}(X) \xrightarrow{\sim} A^{cl}_{\mathfrak{g}}(X)$。

于是为定义（44），只剩构造分次同构 $A^{cl}_{\mathfrak{g}}(X) \xrightarrow{\sim} \mathfrak{z}^{cl}_{{}^{L}\mathfrak{g}}(X)$；这等价于构造概形的 $\mathbb{G}_m$-等变同构 $\Gamma(X,V_{\omega_X}) \xrightarrow{\sim} \operatorname{Hitch}_{{}^{L}\mathfrak{g}}(X)$。由 2.2.2，$\operatorname{Hitch}_{{}^{L}\mathfrak{g}}(X) := \Gamma(X,C_{\omega_X})$，$C := C_{{}^{L}\mathfrak{g}}$。因此只需构造概形的 $\mathbb{G}_m$-等变同构 $V_{\mathfrak{g}} \xrightarrow{\sim} C_{{}^{L}\mathfrak{g}}$。（$V_{\mathfrak{g}}$ 带 3.1.7 中的作用 $a$。）
% [Tutor Note: 原书 p.67 此处与本节末「见 3.1.7」所指的 G_m-作用 a，实际定义于 3.1.9，底本如此，保留原样。]

\SourcePage{68}

由 2.2.1，$C_{{}^{L}\mathfrak{g}} = \operatorname{Spec}(\operatorname{Sym}{}^{L}\mathfrak{g})^{{}^{L}G}$，其中 $G$ 是对应于 $\mathfrak{g}$ 的连通群。可以把 $(\operatorname{Sym}{}^{L}\mathfrak{g})^{{}^{L}G}$ 与 $(\operatorname{Sym}\mathfrak{g}^*)^G$ 等同，因为这两个分次代数都典范同构于 $(\operatorname{Sym}\mathfrak{h}^*)^W$，其中 $W$ 是 Weyl 群。故 $C_{{}^{L}\mathfrak{g}} = C'_{\mathfrak{g}}$，其中
\begin{equation}\tag{45}
C'_{\mathfrak{g}} = \operatorname{Spec}(\operatorname{Sym}\mathfrak{g}^*)^G,
\end{equation}
即 $C'_{\mathfrak{g}}$ 是 $\mathfrak{g}$ 关于 $G$ 伴随作用的仿射概形商。最后，由 Kostant 的工作 [Ko63] 中的定理 0.10，有典范同构 $V_{\mathfrak{g}} \xrightarrow{\sim} C'_{\mathfrak{g}}$，它把 $v \in V_{\mathfrak{g}}$ 映到 $v + i\left(\begin{pmatrix}0&0\\1&0\end{pmatrix}\right) \in \mathfrak{g}$ 在 $C'_{\mathfrak{g}}$ 中的像。该同构与 $\mathbb{G}_m$-作用交换。

\bdp{3.1.14} 描述 $A_{\mathfrak{g}}(X)$ 上的典范滤链与同构（44），还有更自然的方式。

有标准做法把滤 $\mathbb{C}$-代数与分次平坦 $\mathbb{C}[\hbar]$-代数等同（这里 $\deg\hbar = 1$）。具体说，带递增滤链 $\{A_i\}$ 的代数 $A$ 对应分次 $\mathbb{C}[\hbar]$-代数 $A^{\sim} = \bigoplus A_i$，乘以 $\hbar$ 即嵌入 $A_i \hookrightarrow A_{i+1}$。注意 $A = A^{\sim}/(\hbar-1)A^{\sim}$，$\operatorname{gr} A = A^{\sim}/\hbar A^{\sim}$。取谱可知，$\operatorname{Spec} A^{\sim}$ 是直线 $\mathbb{A}^1 = \operatorname{Spec}\mathbb{C}[\hbar]$ 上的平坦仿射概形；$A^{\sim}$ 上的分次等同于 $\operatorname{Spec} A^{\sim}$ 上与 $\mathbb{A}^1$ 的位似作用相容的 $\mathbb{G}_m$-作用。我们要构造概形 $\operatorname{Spec} A_{\mathfrak{g}}(X)^{\sim}$。

设 $\mathfrak{F}$ 是 $X$ 上的 $G$-torsor。记 $\mathcal{E}_{\mathfrak{F}}$ 为 $\mathfrak{F}$ 的无穷小对称构成的李代数胚；有典范正合列
\[
0 \to \mathfrak{g}_{\mathfrak{F}} \to \mathcal{E}_{\mathfrak{F}} \xrightarrow{\pi} \Theta_X \to 0.
\]

回顾：对 $\hbar \in \mathbb{C}$，$\mathfrak{F}$ 上的 $\hbar$-联络是 $\mathcal{O}_X$-线性映射 $\nabla_{\hbar}: \Theta_X \to \mathcal{E}_{\mathfrak{F}}$，满足 $\pi\nabla_{\hbar} = \hbar\operatorname{id}_{\Theta_X}$（通常的联络对应 $\hbar = 1$）。如 3.1.3 那样定义 $G$-$\hbar$-oper，只需把联络 $\nabla$ 换成 $\hbar$-联络 $\nabla_{\hbar}$。关于 $G$-oper 的上述结果都可移用到 $G$-$\hbar$-oper。特别地，$\mathfrak{g}$-$\hbar$-oper（即伴随群的 $\hbar$-oper）构成仿射概形\SourcePage{69}$\operatorname{Op}_{\mathfrak{g},\hbar}(X)$。

对 $\lambda \in \mathbb{C}^*$，有概形同构
\begin{equation}\tag{46}
\operatorname{Op}_{\mathfrak{g},\hbar}(X) \xrightarrow{\sim} \operatorname{Op}_{\mathfrak{g},\lambda\hbar}(X),
\end{equation}
由 $(\mathfrak{F}_B,\nabla_{\hbar}) \mapsto (\mathfrak{F}_B,\lambda\nabla_{\hbar})$ 定义。$\hbar$ 变化时，$\operatorname{Op}_{\mathfrak{g},\hbar}(X)$ 成为仿射 $\mathbb{C}[\hbar]$-概形 $\operatorname{Op}_{\mathfrak{g}}(X)^{\sim} = \operatorname{Spec} A_{\mathfrak{g}}(X)^{\sim}$ 的纤维。对 $\mathfrak{g}$-$\hbar$-oper 使用 3.1.9--3.1.10 的类比，可以证明 $A_{\mathfrak{g}}(X)^{\sim}$ 在 $\mathbb{C}[\hbar]$ 上平坦。态射（46）定义了 $\mathbb{G}_m$ 在 $\operatorname{Op}_{\mathfrak{g}}(X)^{\sim}$ 上的作用，即 $A_{\mathfrak{g}}(X)^{\sim}$ 的分次。$A_{\mathfrak{g}}(X) = A_{\mathfrak{g}}(X)^{\sim}/(\hbar-1)A_{\mathfrak{g}}(X)^{\sim}$ 上相应的滤链与 3.1.13 的滤链一致。

构造（44）等同于构造 $\operatorname{Op}_{\mathfrak{g},0}(X) = \operatorname{Spec}\operatorname{gr} A_{\mathfrak{g}}(X)$ 与 $\operatorname{Hitch}_{{}^{L}\mathfrak{g}}(X) = \operatorname{Spec}\mathfrak{z}^{cl}_{{}^{L}\mathfrak{g}}(X)$ 之间的 $\mathbb{G}_m$-等变同构。如 3.1.11 所述，$\operatorname{Hitch}_{{}^{L}\mathfrak{g}}(X) = \Gamma(X,C'_{\omega_X})$，其中 $C' = C'_{\mathfrak{g}}$ 由（45）定义。有层的典范映射
\begin{equation}\tag{47}
\operatorname{Op}_{\mathfrak{g},0} \longrightarrow C'_{\omega_X},
\end{equation}
它把 $(\mathfrak{F}_B,\nabla_0)$ 映到 $\nabla_0 \in \mathfrak{g}_{\mathfrak{F}}\otimes\omega_X$ 经投影 $\mathfrak{g} \to C'$ 的像。Kostant 的工作 [Ko63] 中的定理 0.10 与命题 19 表明（47）是双射，它诱导所要的同构 $\operatorname{Op}_{\mathfrak{g},0}(X) \xrightarrow{\sim} \Gamma(X,C'_{\omega_X})$。

\bdsec{3.2. 局部 oper 与 Feigin--Frenkel 同构}

\bdp{3.2.1} 把 $X$ 换成形式圆盘 $\operatorname{Spec} O$，$O \simeq \mathbb{C}[[t]]$。3.1 的构造与结果容易移用到这一情形。$\operatorname{Spec} O$ 上的 $\mathfrak{g}$-oper 构成概形 $\operatorname{Op}_{\mathfrak{g}}(O)$，它同构于可数多个变量的多项式环的谱。更确切地说，同构（43）把 $\operatorname{Op}_{\mathfrak{g}}(O)$ 与向量空间 $H^0(\operatorname{Spec} O,V_{\omega_O})$（$V := V_{\mathfrak{g}}$）对应的仿射空间等同。$\operatorname{Spec} O$ 上的 $G$-oper 构成代数栈 $\operatorname{Op}_G(O)$，它同构于 $\operatorname{Op}_{\mathfrak{g}}(O) \times B(Z)$，其中 $B(Z)$ 是中心 $Z \subset G$ 的分类栈，$\mathfrak{g} := \operatorname{Lie}(G/Z)$（该同构不完全是典范的；$\operatorname{Op}_G(O)$ 的典范描述见式（58））。

\SourcePage{70}

与整体情形一样（见 3.1.12--3.1.14），$\operatorname{Op}_{\mathfrak{g}}(O)$ 的坐标环 $A_{\mathfrak{g}}(O)$ 带典范滤链，并有典范同构
\begin{equation}\tag{48}
\sigma_A: \operatorname{gr} A_{\mathfrak{g}}(O) \xrightarrow{\sim} \mathfrak{z}^{cl}_{{}^{L}\mathfrak{g}}(O)
\end{equation}
（见（44））。注意 $\operatorname{Aut} O$ 以显然的方式作用在上述所有对象上。因此 $A_{\mathfrak{g}}(O)$ 是滤 $\operatorname{Aut} O$-代数，$\sigma_A$ 是分次 $\operatorname{Aut} O$-代数的同构。

\bdp{3.2.2} \emph{定理.}（[FF92]）存在滤 $\operatorname{Aut} O$-代数的典范同构
\begin{equation}\tag{49}
\varphi_O: A_{\mathfrak{g}}(O) \xrightarrow{\sim} \mathfrak{z}_{{}^{L}\mathfrak{g}}(O),
\end{equation}
使得 $\sigma_{\mathfrak{z}}\operatorname{gr}\varphi_O = \sigma_A$，其中 $\sigma_{\mathfrak{z}}: \operatorname{gr}\mathfrak{z}_{{}^{L}\mathfrak{g}}(O) \to \mathfrak{z}^{cl}_{{}^{L}\mathfrak{g}}(O)$ 是象征映射。\hfill$\Box$

\emph{注.}
\begin{enumerate}[label=(\roman*),leftmargin=2.2em]
\item 该同构由某些附加的相容性唯一确定；见 3.6.7。
\item Feigin 与 Frenkel 原来的构造是表示论的，对我们来说完全神秘。另一种几何构造见 ???；两种构造的比较见 ???。
\item 对 $\mathfrak{g} = \mathfrak{sl}_2$，（49）有简单的显式描述，本质上属于 Sugawara；见 ???。
\end{enumerate}
% [Tutor Note: 原书 p.70 上述三处交叉引用印作「???」，是 LaTeX 未解析的引用；底本如此，保留原样。]

\bdsec{3.3. 整体版本}

\bdp{3.3.1} 回到整体情形，即 $X$ 是完备曲线。我们将构造代数 $A_{\mathfrak{g}}(X)$ 与 $\mathfrak{z}_{{}^{L}\mathfrak{g}}(X)$ 之间的典范同构（后者由 2.7.4 的公式（27）定义）。
